\section{总结与延伸}

本讲的核心信息是：\textbf{经典 RL 的基本思想（策略梯度、advantage estimation、on/off-policy 区分）在 LLM 推理中同样适用且至关重要}。

关键 takeaway：
\begin{enumerate}
    \item SFT 受限于数据稀缺，纯靠扩展数据的 scaling 速率很慢
    \item RL 通过利用负样本和 advantage 信号，将数据效率提升了约 8 倍
    \item 在线 RL（如 GRPO/PPO）通过持续采样避免了分布偏移
    \item Thinking 模型展示了 RL 可以从简单的 outcome reward 中涌现复杂推理行为
    \item 基础模型的预训练质量是 RL 成功的前提条件
\end{enumerate}

\subsection{方法脉络总表}

\begin{table}[H]
\centering
\begin{tabular}{p{0.18\textwidth} p{0.2\textwidth} p{0.24\textwidth} p{0.22\textwidth}}
\toprule
方法 & 使用的数据/信号 & 主要收益 & 主要风险 \\
\midrule
SFT / NTP & 问题 + 参考解答 & 训练简单、稳定、易扩展 & 受高质量数据规模限制，容易学到表面模式 \\
RFT & on-policy 正样本 + 结果筛选 & 能快速放大已有正确轨迹 & 伪步骤会被一并强化，泛化可能变差 \\
Offline RL & 错误轨迹 + advantage / preference & 能利用负样本，显著提升数据效率 & 依赖 rollout 策略，存在分布偏移 \\
Online RL / GRPO / PPO & 当前策略采样 + outcome / dense reward & 直接优化当前行为分布，适合 thinking model & 成本高，对奖励设计和基础模型质量敏感 \\
\bottomrule
\end{tabular}
\caption{Lecture 10 中几类方法的关系与取舍}
\end{table}

\begin{knowledgebox}{拓展阅读}
\begin{itemize}
    \item DeepSeek-R1: Incentivizing Reasoning Capability in LLMs via RL (2025)
    \item Kimi K1.5: Scaling Reinforcement Learning with LLMs (2025)
    \item Setlur et al., ``RL on Incorrect Synthetic Data Scales the Efficiency of LLM Math Reasoning by Eight-Fold,'' NeurIPS 2024
    \item Gandhi et al., ``Cognitive Behaviors that Enable Self-Improving Reasoners,'' arXiv 2025
    \item Setlur et al., ``Rewarding Progress: Scaling Automated Process Supervision for LLM Reasoning,'' arXiv 2024
\end{itemize}
\end{knowledgebox}

\end{document}
